\section{Decision-Time Counterfactual Supervision}
\label{sec:method}

\Cref{prop:identified} says what a remedy must provide: training samples in which the decision variable moves while
everything else stays, at the states where the policy decides, labeled with the action that is correct after the move.
Decision-time counterfactual supervision (\ours) builds exactly such a $\mathcal{D}_{\text{cf}}$ inside the policy's own
rollouts; each of its components answers one condition of \cref{sec:theory}.

\paragraph{States: the policy's own decisions.}
The student is \pizf with LoRA adapters~\cite{hu2022lora} (rank 32) on the linear layers of the VLM and the action expert;
the teacher is the same checkpoint with the adapters disabled.
Each iteration, the student drives fresh rollouts, and at every queried state $s_t$ the teacher labels the 50-step chunk
$A_t$, as in DAgger~\cite{ross2011dagger}, so that counterfactuals are placed on the states the policy actually visits
(condition (iii)).
The student is trained with the flow-matching loss of \pizf, with $\epsilon\sim\mathcal{N}(0,I)$ and the flow time
$\tau\in(0,1]$ drawn from the Beta-shaped distribution of \pizf ($\tau=1$ is pure noise),
\begin{equation}
\mathcal{L}_{\text{fact}} = \lVert v_\theta(o_t, x_\tau, \tau) - (\epsilon - A_t) \rVert^2,\quad x_\tau = \tau\epsilon + (1-\tau)A_t .
\end{equation}
Without counterfactuals this is our matched control (``Distill''); it leaves \pizf's behavior essentially unchanged, as
\cref{prop:unidentified} predicts.

\paragraph{Interventions: one decision variable at a time.}
At a state $s_t$, the simulator renders a world $s'_t$ in which only the variable that decides the next motion has moved
(condition (ii)):
\begin{itemize}
\item \emph{Retarget} (before the grasp): the target is moved to a free spot 8--30\,cm away, or swaps places with another
object; the hand and the rest of the scene stay.
\item \emph{Relocate} (while carrying): the container, or the furniture region the object goes to, is moved by 15--40\,cm;
hand and object stay.
\item \emph{Co-shift} (before the grasp), a control: hand and target move together, which leaves the decision unchanged and
keeps the factual label correct up to the grasp.
\end{itemize}

\paragraph{Labels: privileged, and correct in the counterfactual world.}
The label of $s'_t$ must be $A(h,x+\Delta)$, not the base policy's output (condition (i)).
We use privileged scripted teachers that know the moved position: an approach to above the target, valid for the steps it
takes to get there (the grasp is left to the policy), and a placer to the new container, valid until shortly after the
release.
A scripted teacher is only as good as its assumptions, and a wrong label identifies a wrong use; three task-agnostic checks
discard the pairs where the scripted notion of the decision variable does not hold:
\begin{itemize}
\item \emph{Uniqueness.} An object or fixture with a twin of the same class is never moved: the instruction can then only
refer to it by where it is (``the bowl next to the plate''), and moving it would change which object is meant.
\item \emph{Agreement.} In the factual world, the scripted teacher must head where the \pizf teacher heads (cosine at least
0.5 between their intended displacements over 25 steps), which rejects states where the task starts with another step or
where the scripted target or container is wrong.
\item \emph{Phase.} Pre-grasp pairs only before the episode's first grasp; relocate pairs only while an object of the task is
held (gripper closed, object lifted and between the fingers).
\end{itemize}

\paragraph{Objective.}
The counterfactual observation $o'_t$ and label $A'_t$, with a validity mask $w$ over chunk steps, enter the same loss with
the same noise and time,
\begin{equation}
\mathcal{L} = \mathcal{L}_{\text{fact}} + \frac{\sum w \odot \lVert v_\theta(o'_t, x'_\tau, \tau) - (\epsilon - A'_t) \rVert^2}{\sum w},
\end{equation}
on a share $\rho=15\%$ of the states of an iteration; \cref{prop:identified} makes $\rho\,\mathbb{E}\lVert\Delta\rVert^2$
the quantity that sets the identification margin. A pairwise consistency term on $v_\theta(o')-v_\theta(o)$ brought no gain
and is not used, in line with~\cite{chen2026ect}.
Objects are recolored at random during the rollouts (probability 0.5 per episode), against a reliance on color that appears
once the policy starts to use the image (\cref{sec:exp:limits}).

\paragraph{Training budget.}
20 iterations of 1024 states each, batch 8, learning rate $5\cdot10^{-5}$ with cosine decay, one suite at a time; every
compared method uses the same teacher, rollouts, budget and optimizer. A run takes 5--7 hours on one H200 GPU, about half
of it simulation and rendering; the counterfactual views add one render per queried state.
Details and the scripted teachers are given in the supplementary material.
